%  LaTeX support: latex@mdpi.com 
%  For support, please attach all files needed for compiling as well as the log file, and specify your operating system, LaTeX version, and LaTeX editor.

%=================================================================
\documentclass[technologies,article,submit,pdftex,moreauthors]{Definitions/mdpi} 

%=================================================================

% MDPI internal commands
\firstpage{1} 
\makeatletter 
\setcounter{page}{\@firstpage} 
\makeatother
\pubvolume{1}
\issuenum{1}
\articlenumber{0}
\pubyear{2023}
\copyrightyear{2023}
%\externaleditor{Academic Editor: Firstname Lastname}
\datereceived{30 January 2023} 
%\daterevised{} % Only for the journal Acoustics
\dateaccepted{} 
\datepublished{} 
\hreflink{https://doi.org/} % If needed use \linebreak
%\doinum{}

%=================================================================
% Add packages and commands here. The following packages are loaded in our class file: fontenc, inputenc, calc, indentfirst, fancyhdr, graphicx, epstopdf, lastpage, ifthen, lineno, float, amsmath, setspace, enumitem, mathpazo, booktabs, titlesec, etoolbox, tabto, xcolor, soul, multirow, microtype, tikz, totcount, changepage, attrib, upgreek, cleveref, amsthm, hyphenat, natbib, hyperref, footmisc, url, geometry, newfloat, caption

\graphicspath{{figures/}} % path to folder with figures
\usepackage{subcaption}
\usepackage{listings}
\usepackage{gensymb} % for \textdegree
\usepackage{tabularx}
\usepackage{spverbatim}

\definecolor{codegreen}{rgb}{0,0.6,0}
\definecolor{codegray}{rgb}{0.5,0.5,0.5}
\definecolor{codepurple}{RGB}{204, 0, 153}
%    
\lstdefinestyle{mystyle}{ 
    numberstyle=\tiny\color{codegray},
    basicstyle=\ttfamily\footnotesize,
    breakatwhitespace=false,         
    breaklines=true,                 
    captionpos=t,                    
    keepspaces=true,                 
    numbers=left,                    
    numbersep=5pt,                  
    showspaces=false,                
    showstringspaces=false,
    showtabs=false,                  
    tabsize=2
}
\lstset{style=mystyle}

%=================================================================
% Full title of the paper (Capitalized)
\Title{GDAL and PROJ libraries integrated with GRASS GIS for terrain modelling \textcolor{blue}{of} the georeferenced raster image}

% MDPI internal command: Title for citation in the left column
\TitleCitation{GDAL and PROJ libraries integrated with GRASS GIS for terrain modelling \textcolor{blue}{of} the georeferenced raster image}

% Author Orchid ID: enter ID or remove command
\newcommand{\orcidauthorA}{0000-0002-5759-1089} % Polina Add \orcidA{} behind the author's name
\newcommand{\orcidauthorB}{0000-0002-6461-1551} % Olivier Add \orcidB{} behind the author's name

% Authors, for the paper (add full first names)
\Author{Polina Lemenkova $^{1}$*\orcidA{} and Olivier Debeir $^{1}$\orcidB{}}

%\longauthorlist{yes}

% MDPI internal command: Authors, for metadata in PDF
\AuthorNames{Polina Lemenkova and Olivier Debeir}

% MDPI internal command: Authors, for citation in the left column
\AuthorCitation{Lemenkova, P.; Debeir, O.}
% If this is a Chicago style journal: Lastname, Firstname, Firstname Lastname, and Firstname Lastname.

% Affiliations / Addresses (Add [1] after \address if there is only one affiliation.)
\address{%
$^{1}$ \quad Laboratory of Image Synthesis and Analysis (LISA), École polytechnique de Bruxelles (Brussels Faculty of Engineering), Université Libre de Bruxelles (ULB). Building L, Campus du Solbosch, ULB -- LISA CP165/57, Avenue Franklin Roosevelt 50, Brussels 1000, Belgium.
}

% Contact information of the corresponding author
\corres{Correspondence: polina.lemenkova@ulb.be; Tel.: +32 471 86 04 59 (P.L.)}

% Current address and/or shared authorship
%\firstnote{Current address: Affiliation 3.} 

% Abstract (Do not insert blank lines, i.e. \\) 
\abstract{Libraries with prewritten codes optimize the workflow in cartography and reduce labour intensive data processing by iteratively applying scripts to implementing mapping tasks. Most existing \textcolor{blue}{Geographic Information System (GIS)} approaches are based on traditional software with a graphical user's interface which significantly limits their performance. Although plugins are proposed to improve the functionality of many GIS \textcolor{blue}{programs}, they are usually ad hoc in finding specific mapping solutions, e.g., cartographic projections and data conversion. We address this limitation by applying the principled approach of \textcolor{blue}{Geospatial Data Abstraction Library (GDAL), library for conversions between cartographic projections PROJ} and \textcolor{blue}{Geographic Resources Analysis Support System (GRASS)} GIS for geospatial data processing and morphometric analysis. This research presents topographic analysis of the dataset using scripting methods which include several tools: 1) GDAL, a translator library for raster and vector geospatial data formats used for converting \textcolor{blue}{Earth Global Relief Model (ETOPO1)} GeoTIFF in XY Cartesian coordinates into \textcolor{blue}{World Geodetic System 1984 (WGS84)} by the ‘gdalwarp’ utility; 2) PROJ projection transformation library used for converting ETOPO1 WGS84 grid to cartographic projections (Cassini-Soldner equirectangular; Equal Area Cylindrical; 2-point Equidistant Azimuthal; Oblique Mercator)\textcolor{blue}{; and }3) GRASS GIS by sequential use of the following modules: r.info, d.mon, d.rast, r.colors, d.rast.leg, d.legend, d.northarrow, d.grid, d.text, g.region, \textcolor{blue}{and} r.contour. The depth frequency was analysed by the module 'd.histogram'. The proposed approach provided a systematic way for morphometric measuring of topographic data and \textcolor{blue}{combine} the advantages of the GDAL, PROJ and GRASS GIS tools that include the informativeness, effectiveness and representativeness in spatial data processing. The morphometric analysis included the computed slope, aspect, profile and tangential curvature of the study area. The data analysis revealed the distribution pattern in topographic data: 24\% of data with elevations below 400 m; 13\% of data with depths -5,000 to -6,000 m; 4\% of depths have values -3,000 to -4,000 m, the least frequent data (-6,000 to 7,000 m) <1\%, 2\% of depths have values -2,000 to 3,000 m in the basin, while other values are distributed proportionally. Further, by incorporating the generic coordinate transformation software library PROJ, \textcolor{blue}{the} raster grid \textcolor{blue}{was transformed} into various cartographic projections to demonstrate distortions in shape and area. Scripting techniques of GRASS GIS are demonstrated for applications in topographic modelling and raster data processing. The GRASS GIS shows the effectiveness for mapping and visualization, compatibility with libraries (GDAL, PROJ), technical flexibility in combining \textcolor{blue}{Graphical User Interface (GUI)} and command-line data processing. The research contributes to the technical cartographic development.}

\keyword{GDAL; PROJ; GRASS GIS; terrain; DEM; elevation; ETOPO1; geomorphometry; cartography; mapping} 

\PACS{\textcolor{blue}{91.10.Da; 92.70.Pq; 93.30.Db; 91.10.Fc}}
\MSC{\textcolor{blue}{86A30; 86A04; 86A22; 53C22; 58E10; 97N60}}
\JEL{\textcolor{blue}{Y91; Q00; Q01; Q5; Q50; Q55; Q56; P18; P48; N50; N57}}

%%%%%%%%%%%%%%%%%%%%%%%%%%%%%%%%%%%%%%%%%%
\begin{document}

\section{Introduction}

\subsection{Background}
A basis form of spatial analysis, topographic modelling consists of visualisation of land surface and plotting contours that reveal variations in elevation \hl{and enables landform delineation} \cite{MINAR2020103414,MAXWELL2022103944}. While in Geographic Information System (GIS) spatial analysis focuses on complex problems regarding the modelling of the Earth's shape \cite{ZHOU2021281,RUZICKOVA2021101109}, technical aspects of spatial data handling consists of various tasks such as data capture \cite{395647,4736372,8647473}, \hl{multi-source data }formatting and processing \hl{ }using functionality of GIS software \cite{7004495,7325718,6946974}, computational analysis and modelling \cite{4659762,5172778,5497752,5653068}, visualization and mapping \cite{938657,6349625,6234685}. Maps became a successful form of geospatial data processing since they help exploiting the complex phenomena of diverse processes and geographic objects on the Earth's surface \cite{PAZOUKI2023108046}. Moreover, maps and\hl{ }georeferenced images enable to learn and recognise regular patterns in visual spatio-temporal datasets for interpretation and analysis of the environmental dynamics and variability \cite{ ,MOREIRA2014208,CHOWDHURY2023102062}.

Therefore, a vital part of the spatial data interpretation is visualization and mapping performed through a cartographic workflow which is a constantly developing process with many existing applications in GIS  \cite{Boulton,Buterez,Duncan}. Due to the variety of the existing GIS (e.g., ArcGIS, QGIS, MapINFO, Environment for Visualizing Images (ENVI), Idrisi, System for Automated Geoscientific Analyses (SAGA) GIS, Erdas Imagine), the choice of the tools may face\hl{ }questions regarding the \hl{technical} functionality of the \hl{diverse} GIS. Among the most essential factors of GIS are the open source availability and functionality providing a variety of tools and embedded computational algorithms \cite{agronomy11050952,Armstrong,ijgi5110209}. Since the \hl{proprietary} GIS such as Erdas Imagine and ArcGIS are a subject of access, open source tools present effective solutions for mapping and spatial data analysis. 

In recent years, open source tools of computer vision and machine learning for data analysis became increasingly successful in analysing spatial data and cartographic tasks \cite{Brovelli}. Free availability of open source GIS cause its popularity and create an increased demand with reported cases in a geospatial analysis \cite{Christl,Turton}. Much progress has been made, for instance, by applications of programming scripts  that aim at automatisation of processing of a given data set. Examples are artificial intelligence and big data processing pointed out in selected studies on geoinformatics \cite{Barnes,Issaka}. The benefits of the open source GIS and scripting are widely discussed and mentioned easily data integration in a number of scenarios and geospatial research, handling 2D and 3D raster datasets \cite{land12010261}, spatial information exchange and processing \cite{Ajvazi,Guan,Kralidis}. 

Open source data, scripting approaches and methods used in geographic research and Earth sciences  \cite{Rey} tend to become progressively popular and perspective in view of the raising amounts of spatial data and the need to process these data accurately yet rapidly. Common sources of open geospatial data include topographic and bathymetric grids \cite{Basith,Habib} using General Bathymetric Chart of the Oceans (GEBCO), Shuttle Radar Topography Mission (SRTM) DEM, ETOPO1 or ETOPO5, geophysical datasets (Earth Gravitational Model of 1996 (EGM96), EGM2008)\hl{, }geological data\hl{, }local data from the online repositories, embedded datasets, OpenStreetMap of Google Earth imagery, remote sensing data such as Landsat TM \cite{YIN2021102457} or Sentinel-2A \cite{app13042525} satellite images.

For many types of spatial data, cartographic modelling approaches hereby have to address the problem of automation while handling these data. An ever-growing pool of spatial datasets encompass multi-format data formats that vary in resolution, source and scale \hl{which raise challenges in handling such big geospatial data} \cite{ZHANG2019101374,KIM2023100368}. These data can be used for the diverse purposes in Earth science such as hydrology and geomorphology \cite{Teixeira}, soils and landscape changes \cite{Savulescu}, environmental coastal studies \cite{Hernandez}, climate changes, land cover studies, geology \cite{app7090956}, and many more. Various GIS software \hl{(e.g., ArcGIS, QGIS, SAGA GIS)} can be applied to process open source geospatial data available online with aim at spatial analysis. Many of the topographic datasets are available in raster of vector digital formats, accessed  via online repositories, e.g., OpenStreetMaps or Google Earth. Increasingly, geospatial data sources are developing online repositories for their data storage such as the National Oceanic and Atmospheric Administration (NOAA), GloVis providing an access to remote sensing datasets. A regular access to such resources enables their use in spatial analysis and arises the question of effective methods for processing these data.

Programming applications address this problem by replicating scripts for treatment of different images and selected scenes of the study area in the topographic map. Other approaches go one step further to automate individual cartographic steps and explicitly model spatial data using modules (as in Generic Mapping Tools -- GMT) or use extensions with prewritten scripts in American Standard Code for Information Interchange (ASCII) format as UNIX shell scripts as in \hl{Geographic Resources Analysis Support System (GRASS) GIS}. Such approaches allow for mapping of datasets by scripts that are independent of their spatial position and replicating these scripts for other extents of data using snippets of code with slightly modified spatial parameters (coordinates, data extent, projections). However, the question of how to better process spatial data and which program to select remains a challenging question in cartography due to many advantages and drawbacks in various tools. 

\subsection{Objectives and Motivation}
This study presents the \hl{the application} of \hl{the }GRASS\hl{ }GIS, an open source powerful software for topographic \hl{modelling}, visualization, statistical\hl{ }analysis and mapping with a case study of\hl{ }morphometric data analysis. Besides open source availability, GRASS GIS includes a variety of functional tools for processing both raster and vector data, using both Graphical User Interface (GUI) and scripts for mapping. Leaving apart the GUI approach as more trivial and easy approach of geospatial data processing, this paper \hl{presents} the \hl{implementation} of scripting methods in \hl{the }GRASS GIS. We presented a case of topographic modelling and data handling using an integration of GRASS GIS with such libraries as Geospatial Data Abstraction Library (GDAL) and PROJ as the key tools for transformation of geospatial data formats and cartographic projections. The combination of these technical tools introduces major computational challenges for cartographic data processing and \hl{selection of the best} projection to avoid distortions \cite{Kiani}.

The GRASS GIS modules were used for assessing the morphometrics in the land-surface parameters over the study area based on the ETOPO1-derived basin morphometry: profile and tangential curvature, slope and aspect maps, histograms of elevation distribution. \hl{To enable data modelling, the QGIS and ArcGIS have a special modelling tool that allows to create new tools with the use of sequential algorithms. However, in contrast to} the traditional GIS software, the architecture of the GRASS GIS presents a collection of modules that can be stored as a sequence in a shell script and used for mapping \hl{using a console}. In this research, we use \hl{such} advantages of the GRASS GIS to demonstrate the effectiveness of scripts for mapping and visualization, as well as its compatibility with libraries, such as GDAL, PROJ and R \cite{Bivand2000,Bivand1999,app122412554,Grohmann}. The theoretical approaches in morphometric studies \cite{Chapman} have been implemented in practical applications of GRASS GIS and developed in GRASS module r.slope.aspect \cite{Hofierka2009}, as demonstrated in this study.  The topographic modelling and spatial analysis based on the continuous field of DEM assist in analysis of slope stability in geologic risk assessment and diverse applications in Earth sciences \cite{HofierkaSuri,Mitas1999,Mitasovaetal1996}. 

\hl{Selected features in the GRASS GIS function comparable to the principle of other GIS. For instance, similar to the QGIS, the }technical flexibility of \hl{the }GRASS GIS enables to switch between the GUI and command-line mode of data processing. Similar to GMT, GRASS GIS uses a programming approach which enables to speed up the repeated tasks through the executable scripts. Like scripts in a UNIX shell, R and Python environments \cite{jimaging8120317}, GRASS GIS and GMT scripts\hl{ }used to automate the workflow repeatedly performed for different maps using the same algorithm for various study areas. Other examples of automation of the cartographic workflow are diverse and include the cases of machine learning approaches of rapid vectorisation of the isolines \cite{Lemenkova2020e}, morphometric landscape analysis \cite{Clarke,Sofia}, geospatial analysis based scripting libraries \cite{su142315966}, topographic mapping using the GMT and \hl{QGIS} \cite{LemenkovaDebeir2022JAES}, fractal modelling quantifying the repetitive patters \cite{Malinverno}.

\section{Case Study}

The study is focused on the Kuril-Kamchatka Trench and Kuril island arc that extend in the north-west Pacific Ocean, from the Kamchatka Peninsula to Hokkaido Island of the Japanese Archipelago\hl{ (Figure }\ref{Fig:1}). This region is characterised by the complex geological settings and high seismicity. \hl{This includes the high discontinuity of the subduction zone caused by the deep slab dynamics and morphology revealed by the seismic cross-sections} \cite{CUI2023117967}. Almost all of islands from the Greater Kuril Chain are located along the submarine ridge of the Kuril island arc with depths mostly 100-300 m above them. The morphological structures include the submerged Vityaz ridge which extends from the oceanward side of the Sea of Okhotsk. The Kuril-Kamchatka Trench generally has a flat seabed consisting of a chain of relatively wide and narrow sections with depths of over -7,000 m. However, it differs in the southern and northern parts with the maximum depth (-9,717 m) in its southwestern segment. 

The basement surface of the Mesozoic structures near the northern coast of the Sea of Okhotsk forms a sub-latitudinal depression with depths reaching up to 3 km. Local minor trenches with depths up to -2.000 m extend from this depression in a sub-meridional direction, covering structural uplift in the centre of the sea. The southern part of the Sea of Okhotsk has a depression with depths exceeding -5,000 m. The steep basement structures of the nearby shelf and the Kuril island arc descend towards the bottom of this depression. 

\vspace{6pt}
\begin{figure}[H]
 \centering
    \includegraphics[width=0.90\textwidth]{Fig_1.jpg}\\
  \caption{Study area on the ETOPO1World Global Relief Model GeoTIFF by NOAA. Mapping: GRASS GIS. Source: authors.}\label{Fig:1}
\end{figure}
\vspace{12pt}

The surface morphology of the seafloor basement in the Kuril-Kamchatka Trench correlates with its modern relief. However, the depths in the seabed of the trench are about -1,000 to -2,000 m deeper. According to the age of the constituent rocks, the Kuril-Kamchatka island arc was formed during the Oligocene period with some of the most ancient seamounts, mostly located on the slopes of the trench, formed on the Mesozoic crust. Following the subduction of the tectonic plates under the island arcs, these seamounts were displaced from the basement surface and imprinted in the structures of the island arcs and coastal structures of active continental margins. Such volcanic mountains are found in the structures of the submarine slopes of the Kamchatka Peninsula. 

The geophysical settings in the study area have a zone of linear positive and negative magnetic anomalies with isometric magnetic field in the Kuril basin, which extends along the Kuril island arc and continues further within the Kamchatka Peninsula. The upper edges of the anomaly-forming bodies, associated with volcanic rocks of the Pliocene-Quaternary age, lie closer to the surface. This briefly demonstrates the existing connectivity and correlation of the submarine and terrestrial morphology which reflects the geological complexity of the region of the Kuril-Kamchatka Trench and Kuril island arc, as briefly discussed above.

%%%%%%%%%%%%%%%%%%%%%%%%%%%%%%%%%%%%%%%%%%
\section{Materials and Methods}

The research was implemented in the GRASS GIS \cite{NetelerMitasova2008}, an open source\hl{ }software and image and graphics production system \cite{Neteleretal2008,Neteler2000}. The major advantage of the GRASS GIS\hl{, similar to the QGIS,} is that it combines the GUI with Command Line Interface (CLI) run from the console\hl{ and enables to operate with cartographic projections by GDAL}. This is convenient to switch between the two modes of data processing since it allows cartographic and analytic experimentation with or without a traditional menu. All maps made in the GRASS GIS were stored using a Universal Transverse Mercator (UTM) map coordinate system \hl{Zone 57} within the geographic extent of the study area. The full scripts used in this study for topographic mapping by the GRASS GIS are available in the GitHub repository of the authors with a provided open access: \href{https://github.com/paulinelemenkova/Mapping_KKT_GRASS_GIS_Scripts}{https://github.com/paulinelemenkova/Mapping\_KKT\_GRASS\_GIS\_Scripts}.

\vspace{6pt}
\begin{figure}[H]
 \centering
    \includegraphics[width=0.95\textwidth]{Fig_2.jpg}\\
  \caption{GDAL library: ‘gdalinfo’ to display raster metadata (left); ‘gdalwarp’ to reproject initial raster file in the XY Cartesian coordinate into WGS 84 (right). Source: authors.}\label{Fig:2}
\end{figure}
\vspace{6pt}

The GRASS GIS uses a modular system organised by the name according to their functionality class (raster, vector) with the first letter indicating a function class, separated by dots and describing the specific task performed by this module. This feature is specific for GRASS GIS, yet differs from the GMT, a scripting cartographic toolset \cite{Wesseletal2019}. Such approach differs both tools from the traditional GIS which rely on the use of GUI. In contrast, the GRASS GIS enables to use both GUI with a standard menu and a scripting approach, which is one of its advantages demonstrating its flexibility between the fully console-based GMT and the conventional GIS. This enables to perform cartographic work in both ways, depending on the needs and preferences and profits from the combination of both approaches. 

\vspace{6pt}
\begin{figure}[H]
 \centering
    \includegraphics[width=0.95\textwidth]{Fig_3.jpg}\\
  \caption{GDAL library: ‘gdal\_translate’ utility to subset target region (left); GRASS GIS command line modules r.info, d.mon, d.rast and visualizing raster (right). Source: authors.}\label{Fig:3}
\end{figure}
\vspace{6pt}

\subsection{Data capture and organisation}

The data used in this project include the open source datasets available online. We used the georeferenced Tag Image File Format (TIFF) file of the ETOPO1 grid by NOAA (\href{https://www.ngdc.noaa.gov/mgg/global/}{NOAA}) as a base map (ETOPO1 Bedrock: Grid of Earth's surface depicting the bedrock underneath the ice sheets), Figure \ref{Fig:1}. The check of metadata for the raster maps was performed by the $‘g.list -f rast’$ GRASS GIS module. The information on the data grids was checked up by command ‘r.info’ module which returns the min/max coordinates of the study area, resolution, number of classes, cartographic projection, scale, geometry with the number of centroids, points, lines, boundaries and islands, and the type of data (3D or 2D). 

\vspace{6pt}
\begin{figure}[h]
 \centering
    \includegraphics[width=0.95\textwidth]{Fig_4.jpg}\\
  \caption{Subset of the study area (Kuril-Kamchatka region), in the Cartesian XY coordinates (metadata: left) visualized on a display (GRASS GIS GUI: right). Source: authors.}\label{Fig:4}
\end{figure}
\vspace{12pt}

The command ‘d.mon’ displays a raster map in a selected monitor. The monitors are used in the GRASS GIS system for visualization of various grids by the commands ‘d.mon wx0’ and ‘d.mon wx1’. 
The project management in GRASS GIS has a certain similarity to the ArcGIS in terms of data storage. Thus, each raster or vector map consists of several auxiliary files which include various data types, categories, header and other specific information. The information on the projection of the current mapset was checked up by the commands: ‘g.proj -p’ (general description), $‘g.proj -w’$ (here: the Area of Interest (AOI) extension and mathematical approximation are as follows: standard parallels, latitude of origin, central meridian, false easting and northing). A summary of the algorithm used for plotting is presented in Listing \ref{lst01}:

\begin{lstlisting}[language=bash,label={lst01},caption=GRASS GIS script-1 (Mapping ETOPO1 raster file)]
#!/bin/sh
# raster file imported to the Location folder 'Kamchatka'
r.in.gdal -o input=/Users/pauline/ETOPO1_KKT_WGS84.tif output=ETOPO1_KKT_WGS84.tif 
# checked up the list of the raster files in the working directory
g.list rast 
#starting graphical display (or 'monitor') in which the maps were displayed
d.mon wx0
# visualizing raster by d.rast utility which displays grid cell maps in the current graphics region
d.rast ETOPO1_Bed_g_geotiff 
#  changing color palette, from default viridis to srtm_plus
r.colors ETOPO1_KKT_WGS84 col=srtm_plus
# check of the range of the elevation values (depths/heights) of the map
r.info ETOPO1_KKT_WGS84
# display a legend on a map
d.rast.leg map= ETOPO1_KKT_WGS84 
# providing numerical output of category numbers or range of values in a raster
r.describe ETOPO1_KKT_WGS84 
#  print a list of category numbers and associated labels
r.cats
# check up and print the current region extent (in this case: NS: 20.000000; EW: 45.000000).
g.region -e 
# Removing auxiliary files
g.remove type=raster name=ETOPO1_KKT_WGS84.tiff -f
\end{lstlisting}

The trial directories were removed from the Unix utility ‘rm’ from a console by commands in two steps: 1) \emph{cd /Users/pauline/grassdata} (changing to the current folder); 2) \emph{rm -r testKKT} (here: removing testing folder for the Kuril-Kamchatka Trench with all files).
The ‘g.region -p’ command was used to retrieve the information on the data extension and selected region (West East South North (WESN), number of rows, columns and cells, aș well aș additional cartographic data: ellipsoid, datum, zone). 

\subsection{Data processing by GDAL}

The initial operations with raster file were performed by the Geospatial Data Abstraction Library (GDAL), a translator library for raster and vector geospatial data formats (GDAL documentation: https://gdal.org/ . The input to the reference list \cite{GDAL} was run from the console of the GMT. In general, the GDAL was used to re-project raster files to the supported cartographic projection, as well as to perform image mosaicing and transformations (Fig. \ref{Fig:2}). In this study, we used three GDAL utilities using the commands of the following utilities:\\ 
$‘gdalwarp’, ‘gdal\_translate’ and ‘gdalinfo’$ (Fig. \ref{Fig:2}). The ‘gdalwarp’ utility was used to re-projection and warping of the initial original file (ETOPO1\_Bed\_g\_geotiff.tif) from the XY Cartesian coordinates to the World Geodetic System 1984 (WGS84) \cite{Defense}. 

The transformation uses the European Petroleum Survey Group (EPSG) number, which in this case is 4326 (corresponds to the WGS84). The warping was done by the command for Spatial Reference System (SRS) as follows: 

\begin{spverbatim}
gdalwarp -t_srs EPSG:4326 ETOPO1_Bed_g_geotiff.tif ETOPO1_WGS84.tif
\end{spverbatim} 

The next step contains running the utility ‘gdal\_translate’  (Fig. \ref{Fig:3}) with the argument $‘-projwin\_srs <srs\_def>’$, which specifies the projwin SRS GDAL, in which it interprets the coordinates (here, given coordinates in the EPSG system). 

Using the gdal\_translate (Figure \ref{Fig:3}), the file was subset, or clipped, to the target square area from the global extent with clockwise going coordinates in degrees as follows: W=140 N=60 E=170 S=40. The re-projection was done by the ‘grdalwarp’ utility. The subsetting of the target area was done by the following command: 

\begin{spverbatim}
gdal_translate -projwin 140.0000 60.0000 170.0000 40.0000 ETOPO1_WGS84.tif 
ETOPO1_KKT_WGS84.tif
\end{spverbatim}

\vspace{6pt}
\begin{figure}[h]
 \centering
    \includegraphics[width=0.95\textwidth]{Fig_5.jpg}
  \caption{Raster projected from XY to WGS84 Lat/Lon coordinates with cartographic elements. Source: authors.}\label{Fig:5}
\end{figure}
\vspace{12pt}

Now that all the components of the code have been explained, we provide a summary in Listing \ref{lst02}
\begin{lstlisting}[language=bash,label={lst02},caption=GRASS GIS script-2 (adding cartographic elements)]
#!/bin/sh
g.region rast=ETOPO1_KKT_WGS84.tif -p
d.mon wx0
d.rast ETOPO1_KKT_WGS84
r.colors ETOPO1_KKT_WGS84 col=srtm_plus
d.legend raster=ETOPO1_KKT_WGS84 title=Elevation,m title_fontsize=12 \
	font=Helvetica fontsize=8 -t -b bgcolor=white label_step=2000 border_color=gray -f thin=10
d.northarrow style=fancy_compass rotation=0 label=N -w color=black fill_color=gray fontsize=10
d.grid size=02:00:00 -a -d color=red width=0.1 fontsize=8 text_color=white
d.text text="Region of Kuril-Kamchatka Trench" color=yellow bgcolor=gray size=8
\end{lstlisting}

The method presentation is shown by the projection transformation by GDAL utility ‘gralwarp’ using the PROJ library. It demonstrates the transformation of the TIFF file in the WGS84 geoid datum warped to a Two Point Equidistant Azimuthal projection using the PROJ arguments: 

\begin{spverbatim}
gdalwarp -t_srs '+proj=tpeqd +lat_1=45 +lat_2=55 +lon_1=100 +lon_2=120' 
ETOPO1_KKT_WGS84.tif ETOPO1_KKT_2p.tif -overwrite
\end{spverbatim}

Now the file was reprojected to the WGS84 and subset (cut off) from the initial raster grid with global extent. A gdalinfo utility was used to check up the correctness of the performed operations and the output data as follows: $gdalinfo ETOPO1\_KKT\_WGS84.tif.$ Next, we present the code in details in Listing \ref{lst03}:

\begin{lstlisting}[language=bash,label={lst03},caption=GRASS GIS script for plotting 2 Point Equidistant Azimuthal projection]
#!/bin/sh
g.list rast
r.info ETOPO1_KKT_2p
d.mon wx0
g.region raster=ETOPO1_KKT_2p -p
r.colors ETOPO1_KKT_2p col=grey.eq
d.rast ETOPO1_KKT_2p
d.redraw
r.contour ETOPO1_KKT_2p out=Bathymetry1000 step=1000
d.vect Bathymetry1000 color='blue' width=0
d.grid -g size=2.5 color='255:0:0'
d.text text="2 Point Equidist. Azimuth." color='0:0:51' bgcolor='224:224:224' size=3
\end{lstlisting}

The next step was continued in the GRASS GIS version 7.6 and included the coordinates transformation that was done by ‘proj’ with several tested projections and their visualization using GRASS GIS scripts presented in details in Listings \ref{lst04} for Cassini-Soldner projection, Listing \ref{lst05} for Oblique Mercator projection and Listing \ref{lst06} for and Equal-Area Cylindrical projection. The ‘proj’ is a cartographic projection filter of the PROJ software which allows converting between the geographic and projection coordinates within one datum WGS84. The maps projections were tested in PROJ via the ‘gdalwarp’ utility of GDAL before importing them into the GRASS GIS. The ‘gdalwarp’ is specially designed to reproject maps to the target projections, coordinate systems, ellipsoid, or geodetic datum. In this case, we selected the datum as a constant WGS84 and kept it constant for this project. 

\begin{lstlisting}[language=bash,label={lst04},caption=GRASS GIS script for plotting map in Cassini-Soldner projection)]
#!/bin/sh
g.list rast
r.info ETOPO1_KKT_C
d.mon wx0
g.region raster=ETOPO1_KKT_C -p
r.colors ETOPO1_KKT_C col=grey.eq
d.rast ETOPO1_KKT_C
d.redraw
r.contour ETOPO1_KKT_C out=Bathymetry1000C step=1000
d.vect Bathymetry1000C color='blue' width=0
d.grid -g size=2.5 color='255:0:0'
d.text text="Cassini-Soldner prj" color='0:0:51' bgcolor='224:224:224' size=3
\end{lstlisting}

\begin{lstlisting}[language=bash,label={lst05},caption=GRASS GIS script for plotting map in Oblique Mercator projection)]
#!/bin/sh
g.list rast
r.info ETOPO1_KKT_OM
d.mon wx0
g.region raster=ETOPO1_KKT_OM -p
r.colors ETOPO1_KKT_OM col=grey.eq
d.rast ETOPO1_KKT_OM
r.contour ETOPO1_KKT_OM out=Bathymetry1000OM step=1000
d.vect Bathymetry1000OM color='blue' width=0
d.grid -g size=2.5 color='255:0:0'
d.text text="Oblique Mercator prj" color='0:0:51' bgcolor='224:224:224' size=3
\end{lstlisting}

\begin{lstlisting}[language=bash,label={lst06},caption=GRASS GIS script for plotting map in Equal-Area Cylindrical projection)]
#!/bin/sh
g.list rast
r.info ETOPO1_KKT_EAC
d.mon wx0
g.region raster=ETOPO1_KKT_EAC -p
r.colors ETOPO1_KKT_EAC col=grey.eq
d.rast ETOPO1_KKT_EAC
d.redraw
r.contour ETOPO1_KKT_EAC out=Bathymetry1000 step=1000
d.vect Bathymetry1000 color='blue' width=0
d.grid -g size=2.5 color='255:0:0'
d.text text="Equal Area Cylindrical prj" color='0:0:51' bgcolor='224:224:224' size=3
\end{lstlisting}

\subsection{Displaying raster grids in GRASS GIS}

Upon re-projection of the file, it was imported to the GRASS GIS by the ‘r.in.gdal’ utility. Displaying the dataset was done using the sequence of the GRASS modules shown in Listing 1, based on the GRASS GIS modular techniques \cite{Brown} and visualized in Figure \ref{Fig:4}. The visualized cartographic elements include grid ticks, legend and the title, which were placed on the map in Figure \ref{Fig:5} using the code in Listing 2. Generating the continuous fields of the bathymetric contours (isolines) were derived from the raster ETOPO1\_KKT\_WGS84 file using the GRASS GIS module 'r.contour' by script. In this example the step parameter was set to 750 m and visualized in Figure \ref{Fig:6}. The script of the GRASS GIS used for plotting the topographic isolines is shown in Listing \ref{lst07}:

\begin{lstlisting}[language=bash,label={lst07},caption=GRASS GIS script for visualising raster grid and adding contour isolines]
#!/bin/sh
d.mon wx0
r.info ETOPO1_KKT_WGS84
g.region raster=ETOPO1_KKT_WGS84 -p
d.rast ETOPO1_KKT_WGS84
d.grid size=03:00:00 -a -d color=red width=0.1 fontsize=8 text_color=white
d.text text="Kuril-Kamchatka Area" color='255:255:255' size=5
r.contour ETOPO1_KKT_WGS84 out=Bathymetry750 step=750 --overwrite
d.vect Bathymetry750 color='102:37:4' width=0
\end{lstlisting}

\vspace{6pt}
\begin{figure}[h]
 \centering
    \includegraphics[width=0.95\textwidth]{Fig_6.jpg}\\
  \caption{Plotting isolines by the 'r.contour' module and visualising raster in ‘elevation’ color table. Source: authors.}\label{Fig:6}
\end{figure}
\vspace{12pt}

%%%%%%%%%%%%%%%%%%%%%%%%%%%%%%%%%%%%%%%%%%
\section{Results and Discussion}
We implemented the GRASS GIS algorithms for cartographic data visualization and compiled maps in various projections. The topographic data visualization and a morphometric analysis were based using the modules of the GRASS GIS. The input cartographic data included the ETOPO1 raster grid including the derived variables: slope, aspect, curvature and a hillshade. The data analysis included a plotted histogram and a pie chart showing data distribution. Several projections were tested and visualized using PROJ library. The data transformation was performed using GDAL integrated with the GRASS GIS. 

\subsection{Transforming raster grid into cartographic projections by PROJ}

We used PROJ, a generic coordinate transformation software library (https://proj.org/), to transform the raster grid files in Lat/Lon geospatial coordinates from one coordinate reference system (CRS) to another. This method includes the PROJ transformation which performs a variety of operations with cartographic data, starting from the simple and well-known projections such as cylindrical Mercator projection to the very complex transformations across many reference coordinate frames. The PROJ was used for testing various projections (conic, cylindtrical, azimuthal) selected from the variety of available projections in the library: Cassini-Soldner equirectangular; Equal Area Cylindrical; 2-point Equidistant Azimuthal; Oblique Mercator \cite{Snyder}. We installed the PROJ library via the Homebrew package manager using the following commands: ‘brew install proj’ and ‘brew upgrade proj’ for installing and updating the package, respectively. 

Originally made as a technical tool for cartographic projections, PROJ gradually transformed into being a functional generic coordinate transformation engine. The supported projections were checked up by the command ‘cs2cs -lp’ from the GRASS GIS console \cite{Evenden}. The PROJ was applied for the scale cartographic projections and coordinate transformation (Figure \ref{Fig:7}), thus presenting the geodetic transformations of varying complexity. In this study, the visualization of various cartographic projections in GRASS GIS (Figure \ref{Fig:7}) was performed by a combination and sequential use of the following modules of GRASS GIS: 'd.rast', 'r.contour', 'd.vect', 'd.grid' and 'd.text'. The script is based on the developed methodology of data processing by GRASS GIS \cite{Dassau}. It was also used for plotting the contour maps with an example of Cassini-Soldner projection. Mapping study area in various cartographic projections was recorded in separate plots in the GRASS GIS and  presented in Figure \ref{Fig:7} for comparison: Cassini-Soldner equirectangular; Equal Area Cylindrical; Two Point Equidistant Azimuthal; Oblique Mercator. 
 
 \vspace{6pt}
\begin{figure}[h]
 \centering
    \includegraphics[width=0.90\textwidth]{Fig_7.jpg}
  \caption{Comparison of the four maps visualized in GRASS GIS in four different cartographic projections transformed by the ‘gdalwarp’ GDAL utility and PROJ parameters. The projections from left to right, top to bottom are as follows: Cassini-Soldner equirectangular; Equal Area Cylindrical; Two Point Equidistant Azimuthal; Oblique Mercator. Source: authors.}\label{Fig:7}
\end{figure}
\vspace{12pt}

\subsection{Calculating morphometric parameters}

The correlation between the profile and tangential curvature reflects the methodological approaches of the algorithm which concerns the mathematical approximation of the morphological profile of the slopes and the inflation of the terrain heights. The differences in calculus are well illustrated in the Equation \ref{pc} for the profile curvature $K_P$ \cite{Golden} and the Equation \ref{tc} for the tangential curvature $K_T$ \cite{Golden}. These formulae are given for the profile curvature in Equation \ref{pc} and for the tangential curvature in Equation \ref{tc}, respectively. 

\begin{equation}
  \label{pc}
 K_P=\frac{\left(\frac{\delta^2 z}{\delta x^2}\right) \left(\frac{\delta z}{\delta x}\right)^2 + 2\left(\frac{\delta^2 z}{\delta x \delta y}\frac{\delta z}{\delta x}\frac{\delta z}{\delta y} \right) + \frac{\delta^2 z}{\delta y^2} \left(\frac{\delta z}{\delta x}\right)^2}{pq^\frac{3}{2}}
\end{equation}

\begin{equation}
  \label{tc}
 K_T=\frac{\left(\frac{\delta^2 z}{\delta x^2}\right) \left(\frac{\delta z}{\delta y}\right)^2 - 2\left(\frac{\delta^2 z}{\delta x \delta y}\frac{\delta z}{\delta x}\frac{\delta z}{\delta y} \right) + \frac{\delta^2 z}{\delta y^2} \left(\frac{\delta z}{\delta x}\right)^2}{pq^\frac{1}{2}}
\end{equation}

where $p=\left(\frac{\delta z}{\delta x}\right)^2+\left(\frac{\delta z}{\delta y}\right)^2$ and $q=1+p$ for both cases.
 
The input elevation data were derived from the ETOPO1 raster grid that contained the elevation values in digital format covering study area. The 'r.slope.aspect' module of GRASS GIS was used for plotting raster maps using the developed methodology \cite{Mitasova1985,Hofierka2009,Horn}. The resulting maps of slope, aspect (computed counterclockwise from the east), profile curvature and tangential curvature from the ETOPO1 raster grid are shown in Figure \ref{Fig:9}.

The plotting was done by the code: 

\begin{spverbatim}
‘r.slope.aspect elevation=ETOPO1_KKT_WGS84 slope=slope aspect=aspect 
pcurvature=pcurv tcurvature=tcurv’
\end{spverbatim}

where the input elevation raster map is the ETOPO1\_KKT\_WGS84, and the output raster map is ‘slope’. The results include the morphometric elements of slope, aspect, profile curvature and tangential curvature, Figure \ref{Fig:9}.

\vspace{6pt}
\begin{figure}[H]
 \centering
    \includegraphics[width=0.90\textwidth]{Fig_9.jpg}
  \caption{Slope, aspect, profile and tangential curvature maps, Kuril-Kamchatka region, based on raster grid ETOPO1. Source: authors.}\label{Fig:9}
\end{figure}
\vspace{12pt}

The computed geomorphic indices included slope, aspect, profile and tangential curvature maps, derived from the topographic grid using the embedded algorithms in GRASS GIS. The curvature shows the rate of variations of the slope, which makes it the second function of the elevation variable related to the surface heights \cite{Schmidt}. This process is technically both both cases: whether the purpose is profile curvature or tangential curvature. However, the difference is in the approaches of the mathematical algorithms that are embedded in the GRASS GIS. Hence, the difference between the profile and tangential curvature consists in the method of calculation. For the tangential curvature, it is computed as a perpendicular to the slope gradient, while for the profile curvature, it measures a curvature of the surface in the direction of the gradient \cite{Gallant}. 

The slope shows a steepness or the degree in the topographic surface covering the study area. The aspect, or exposure, is a compass direction that a slope faces (East, West, South, North). Geographically, the aspect of a slope has a significant impact on the local microclimate which affects physical and often botanical features of a mountain slope, i.e., a slope effect on the terrestrial areas. Such effects are owing to the strong correlation between the mountain slope and the temperature, which is in turn affected by the angle of sun. Hence, both parameters are derivative from the slope of the elevation with certain variations in the mathematical approach and angle direction of the gradient degree (direct and perpendicular). The theoretical foundations of the slope and angle modelling in geospatial research can further be found in the existing works reporting, for instance, on the effects from the cell size \cite{Hodgson1995}, gradient angles \cite{Hodgson1998}, numerical terrain analysis \cite{Zevenbegen}, or quantifying land surface and slope curvatures \cite{Dikau,Dunn,Guth,Heerdegan}. 

\subsection{Analysis of elevation frequency}
The next research step included the analysis of the distribution of the elevation over the study area. Using computational techniques of GRASS GIS, we calculated the frequency of the data in various range diapason and repeatability of corresponding depths. The most frequent depths or elevations over the study area were recorded and summarised for the analysis of spatially distributed data \cite{Antrop,Bailey,Burrough}. The increase in values of a specific depth reflects the frequency of the elevation recorded on a bathymetric grid in a subsequent region of the study area. This was performed by plotting the histograms as bars and pie charts using the GRASS GIS 'd.histogram' module for calculating depths distribution \cite{Mitasovaetal1995}. 

The analysis of the spatial data distribution aimed to demonstrate a range of variations in data, including skewness, the variation of values within intervals of the topography, organisation of categorical data, selecting relevant peaks in the interpretation and description of the dataset or other resources showing the correlation among the geographic phenomena such as aspect and slope with relief. The distribution of topographic data has been evaluated using statistical tools embedded in GRASS GIS and presented in a histograms and a pie chart (Figure \ref{Fig:8}).

\vspace{6pt}
\begin{figure}[h]
 \centering
    \includegraphics[width=0.90\textwidth]{Fig_8.jpg}\\
  \caption{Bar graph histogram for elevation map (left), pie graph of heights distribution based on the raster grid ETOPO1. Source: authors.}\label{Fig:8}
\end{figure}
\vspace{12pt}

Critical synthesis and analysis of a geomorphological setting of the study area is largely based on the experiential techniques of cartographic data visualization using GRASS GIS. A range of its modules utilised in scripts for cartographic transformations were used to visualise various projection and colour map palettes. The topographic data analysis demonstrated variations in data distribution and morphometry in the study area, as follows. The percentage of the topographic data distribution over the study area is as follows. The depths in the interval below 400 m to -2,000 m covered 24\% of the total points; the depths in the interval from -2,000 to -3,000 m sum up to 2 \%; the depth  from -3,000 to -4,000 m take up 4\% of the total points; the depths in the range between -5,000 to -6,000 m are 13\%, and the depth in the diapason from -6,000 to 7,000 are less than 1\%. the rest of the elevation points is covered by the data with diverse heights and topographic values.
  
The extent of the data range is summarised and visualized in the histogram and a pie chart, Figure \ref{Fig:8}). The bar plot and the pie chart show the frequencies of the elevation levels in a raster representing morphometric unevenness of the terrain. For each topographic range, the number of pixels in the raster image is counted and drawn on the diagram in the two views of representation. The number of pixels corresponding to the each value of depth/height is contributing to the height of this particular bar on the chart (or a radius of the segment in a pie). The x-axis of the bar chart diagram represents the topographic/bathymetric values (heights and depths), while vertical y-axis shows the number of pixels counted at a given step.

The results of the morphometric modelling show that the terrain slope in the Kuril-Kamchatka region is ranging from 0° to 39°, as shown on the slope map calculated from the ETOPO1\_KKT\_WGS84 grid. The slope in the study area is shown in gradually increasing degrees starting from 0° (flat horizontal) to the 39° (maximal vertical). The profile curvature shows modelled terrain surface in the gradient direction which reflects the variations in the slope steepness angle within the range of -87*10-5 to 68*10-5. The most part of the area (coloured light green on the lower left map, Figure \ref{Fig:9}) has values of -10*10-5 to 0. Practically, it indicates that the velocity of mass falling downwards along the slope curve. The tangential curvature reflects the change in the aspect angle and influences the divergence, or the convergence, of the water flow, respectively.

The tangential curvature shows values of -67*10-5 to 63*10-5. The terrestrial areas of the Kamchatka Peninsula and Greater Kuril Chain, respectively, show higher positive values (depicted by the dark grey colours, in the lower right map in Figure \ref{Fig:9}) comparing to the water areas of the Sea of Okhotsk (white and light grey colours) which shows variations in the land surface steepness over the terrain compared to water areas. Numerically, the aspect map shows the orientation in degrees varying from 0° to 360° (upper right map in Figure \ref{Fig:9}). It shows a clear correlation between the aspect and the topographic area of the Kuril-Kamchatka Trench (130°-150°, emerald green colour on the map) and the Greater Kuril Chain which extends in parallel to the trench both in south-eastern orientation (beige colour from the selected ‘aspectcolr’ colour palette, 300°-330°). A clear delineation of the Sakhalin Island can be noted as the north-western orientation (yellow-green colours, 0°-30°) and the Kamchatka Peninsula (NW, greens, 60°-120°), respectively.


%%%%%%%%%%%%%%%%%%%%%%%%%%%%%%%%%%%%%%%%%%
\section{Discussion}

The visualized images are displayed using various colour map palettes (r.colors) specifically intended for the cartographic visualization of the terrain raster grids. The produced maps show morphometric variables over the target area performed by the GRASS GIS techniques which included data conversion, warping, transforming projections and adding cartographic elements (ticks, legend, isolines) on the final layout. Each map sources a range of new visualization of the data relevant to the morphology of the study region. Similar to the GMT techniques, the scripting methods  of the GRASS GIS have a command-line approach of data processing. This is convenient for controlling the cartographic design and enables to create the print-quality maps. Further, we demonstrated the effectiveness of the GRASS GIS methods in visualising raster grids and performing raster data analysis in a geologically complex regions such as the north-western sector of the Pacific Ocean.

The morphometry is an essential approach of the quantitative GIS-based terrain analysis. It is increasingly used in diverse applied geological tasks for relief visualization. Based on elevation of DEM, the landforms can be visualised in the target area for morphological modelling and mapping at regional scale. The morphometric terrain analysis may further include extended approaches such as analysis and interpretation of the land forms or seafloor surface as a continuous field, upscaling the morphological shapes for a closer analysis of the terrain structure. The slope gradient indicates the terrain roughness and represents a basic for the land surface discreteness. A slope gradient can be used for indication of the terrain roughness, while the aspect can be applied for mapping slope distribution, depending on the compass orientation of the terrain. 

The findings revealed that the topography of the Kuril-Kamchatka Trench varied significantly reflecting its regional geologic evolution and tectonic setting which largely affects the structure of the relief. Our study presented a case for displaying topographic data for analysis of height differences using the cartographic and statistical data processing. We shown the example of bathymetric data processing specifically for the north Pacific Ocean area performed using scripting approach by GRASS GIS syntax which has its specific language operated from the command line. The scripts of the GRASS GIS were used to rapidly plot the topographic data as a series of morphometric variables. The results can be extended to similar datasets using scripts presented for repeatability of data processing in GRASS GIS: visualization of the slope, aspect, curvature and statistical data distribution of elevation data.

%%%%%%%%%%%%%%%%%%%%%%%%%%%%%%%%%%%%%%%%%%
\section{Conclusions}

We presented an application of the GRASS GIS algorithms for topographic mapping and morphometric modelling that uses a script-based approach for cartographic data visualization, transformation and analysis. The selected region of the tested area includes the north-western segment of the Pacific Ocean in the Far East. Extensive experimental results on mapping raster images by scripts taken from the ETOPO1 grid show that GRASS GIS outperforms the traditional software in terms of automation of data processing which increase the speed and accuracy of cartographic plotting. We have also shown that the geodetic transformation of the dataset in various cartographic projections performs well using GDAL utility and PROJ library. 

While the choice of the cartographic projections is wholly due to the spatial location of the study area controlled by the possible distortions in sizes, angles, distances and directions of the visualized objects, we demonstrated the method of flexible changing of projections using the integration of GRASS GIS and GDAL. With regard to addressing the cartographic plotting, GRASS GIS show robustness and effectiveness of data handling both using the GUI interface and from the console in a series of presented listings. We believe that techniques of GRASS GIS for handling geospatial data can be further extended to deal with spatial datasets using scripts presented in this study.

The geomorphometric parameters include slope, aspect, and curvature in both the down‐slope and across‐slope directions, which is presented in this case study based on the GRASS GIS modules. The results show the effective application of the GRASS GIS techniques for cartographic data processing. Using the console by scripting mode, we identified topographic variability in the target area through spatial analysis of the processed data and visualised raster grids in various cartographic projections. 

%%%%%%%%%%%%%%%%%%%%%%%%%%%%%%%%%%%%%%%%%%
\authorcontributions{Supervision, conceptualization, methodology, software, resources, funding acquisition, and project administration, O.D.; writing---original draft preparation, methodology, software, data curation, visualization, formal analysis, validation, writing---review and editing, and investigation, P.L. All authors have read and agreed to the published version of the manuscript.}

\funding{The publication was funded by the Editorial Office of Technologies, Multidisciplinary Digital Publishing Institute (MDPI), by providing 100\% discount for the APC of this manuscript. This project was supported by the Federal Public Planning Service Science Policy or Belgian Science Policy Office, Federal Science Policy - BELSPO (B2/202/P2/SEISMOSTORM).}

\institutionalreview{Not applicable.}

\informedconsent{Not applicable.}

\dataavailability{Not applicable} 

\acknowledgments{The authors thank the anonymous reviewers of Land for reading, suggestions and comments that improved an earlier version of this manuscript.}

\conflictsofinterest{The authors declare no conflict of interest.} 

\abbreviations{Abbreviations}{
The following abbreviations are used in this manuscript:\\

\noindent 
\begin{tabular}{@{}ll}
ASCII & American Standard Code for Information Interchange \\
AOI & Area of Interest \\
CLI & Command Line Interface \\
DEM & Digital Elevation Model \\
EGM96 & Earth Gravitational Model of 1996 \\
EGM2008 & Earth Gravitational Model of 2008 \\
ENVI & Environment for Visualizing Images \\
EPSG & European Petroleum Survey Group \\
ETOPO1 & Earth's surface topography \\
GDAL & Geospatial Data Abstraction Library \\
GEBCO & General Bathymetric Chart of the Oceans\\
GIS & Geographic Information System \\
GMT & Generic Mapping Tools \\
GRASS & Geographic Resources Analysis Support System \\
GUI & Graphical User Interface \\
NOAA & National Oceanic and Atmospheric Administration \\
PROJ & a library for performing conversions between cartographic projections \\
QGIS & Quantum GIS \\ 
SAGA & System for Automated Geoscientific Analyses \\
SRS & Spatial Reference System \\
SRTM & Shuttle Radar Topography Mission\\
TIFF & Tag Image File Format \\
UTM & Universal Transverse Mercator \\
WGS84 & World Geodetic System 1984\\
WESN & West East South North
\end{tabular}
}

%%%%%%%%%%%%%%%%%%%%%%%%%%%%%%%%%%%%%%%%%%
\begin{adjustwidth}{-\extralength}{0cm}
%\printendnotes[custom] % Un-comment to print a list of endnotes

\reftitle{References}

%=====================================
% References, variant A: external bibliography
%=====================================
\bibliography{TechBib.bib}

%%%%%%%%%%%%%%%%%%%%%%%%%%%%%%%%%%%%%%%%%%
\end{adjustwidth}
\end{document}
